\documentclass[11pt]{article}
\usepackage{docstyle}
\renewcommand\sheetname{Handout 4 \textperiodcentered{} Merit and Excellence problem solving}

\begin{document}
\sheethead{Lesson 4 \textperiodcentered{} Merit and Excellence Problem Solving}
  {NCEA Level 2 Mathematics \textperiodcentered{} Algebra}{Handout}

\begin{keybox}[title={The one idea: general, complete, checked}]
At Excellence the algebra is the same as in Lessons 1--3. What changes is the \textbf{argument}: it must hold for \textbf{every} case (letters, not examples), find \textbf{every} answer, and be \textbf{checked} against the question, with a written conclusion. Markers write it as ``correct mathematical statements'' and ``with reasoning''.
\end{keybox}

% =================================================================== 1
\section{Algebraic proof}\label{h4:proof}
\begin{minipage}[t]{0.46\textwidth}\vspace{0pt}
\renewcommand\arraystretch{1.4}
\begin{tabular}{@{}>{\raggedright}p{0.52\linewidth} >{\raggedright\arraybackslash}p{0.4\linewidth}@{}}
\toprule
\textbf{Words} & \textbf{Algebra}\\ \midrule
consecutive whole numbers & $n$, $n+1$, $n+2$ \\
an even number & $2n$ \\
an odd number & $2n+1$ \\
two different odd numbers & $2m+1$, $2n+1$ \\
a multiple of 3 & $3k$ \\
next day down a calendar & $n+7$ \\
\bottomrule
\end{tabular}\par\smallskip
$n$, $m$, $k$ stand for whole numbers.
\end{minipage}\hfill
\begin{minipage}[t]{0.5\textwidth}\vspace{0pt}
\textbf{Write it in this order}
\begin{enumerate}
\item Name every number with letters.
\item Write the expression the question is about, and simplify it.
\item Take out the factor the question needs: $3(n+1)$, $8n$, $2(\ldots) + 1$.
\item Finish with a sentence that says what the factor means: ``so it is always a multiple of 3''.
\end{enumerate}
Examples, however many, are not a proof.
\end{minipage}

\begin{multicols}{2}
\textbf{Example 1.} Prove that the sum of three consecutive whole numbers is a multiple of 3.
\[ n + (n+1) + (n+2) = 3n + 3 = 3(n+1) \]
$3(n+1)$ is 3 times a whole number, \textbf{so the sum is always a multiple of 3.} \mtag{r}

\textbf{Example 2.} Prove that $(2n+1)^{2} - (2n-1)^{2}$ is a multiple of 8.
\[ (4n^{2}+4n+1) - (4n^{2}-4n+1) = 8n \]
\textbf{so it is always a multiple of 8.} \mtag{r}
\columnbreak

\textbf{Example 3.} On a calendar, a 3-by-3 square has top-left number $n$. Show that the diagonals have equal sums but never equal products.\par
\par\smallskip{\setlength\figunit{4mm}\hbox to\linewidth{\hfil\FigCalAlgebra\hfil}}\par\smallskip
Corners: $n$, $n+2$, $n+14$, $n+16$.\par
Sums: $n + (n+16) = 2n + 16 = (n+2) + (n+14)$.\par
Products: $(n+2)(n+14) - n(n+16) = 28$, never 0, \textbf{so the products are never equal.} \mtag{t}
\end{multicols}

% =================================================================== 2
\section{Where a line meets a curve}\label{h4:simul}
\begin{keybox}[title={Write it in this order}]
\begin{enumerate}
\item Make one variable the subject of the simpler equation ($y = \ldots$ or $x^{2} = \ldots$).
\item Substitute into the other: you get a quadratic in one variable. Rearrange to $= 0$.
\item Solve, then find the \textbf{other coordinate for every value}. Reject values that give no real point.
\item Tangent (touches once): the quadratic has $b^{2} - 4ac = 0$. Do not meet: $< 0$.
\end{enumerate}
\end{keybox}

\begin{multicols}{2}
\textbf{Example 4.} Where does $y = 2x - 3$ meet $y = x^{2} - 2x - 3$?
\begin{align*}
x^{2} - 2x - 3 &= 2x - 3\\
x^{2} - 4x = x(x - 4) &= 0
\end{align*}
$x = 0$: $y = -3$;\quad $x = 4$: $y = 5$.\quad \textbf{$(0,-3)$ and $(4,5)$} \mtag{u}

\textbf{Example 5.} Find $k$ so that $y = kx$ is a tangent to $y = x^{2} + 3$.
\[ x^{2} - kx + 3 = 0,\quad k^{2} - 12 = 0,\quad \mathbf{k = \pm 2\sqrt{3}} \]
\mtag{r}
\columnbreak

\textbf{Example 6.} Where do $x^{2} + y^{2} = 13$ and $y = x^{2} - 1$ meet?\par
$x^{2} = y + 1$, so $y + 1 + y^{2} = 13$:
\begin{align*}
y^{2} + y - 12 = (y+4)(y-3) &= 0
\end{align*}
$y = 3$: $x^{2} = 4$, $x = \pm 2$.\quad $y = -4$: $x^{2} = -3$, \textbf{no real point}.\par
Points: \textbf{$(2, 3)$ and $(-2, 3)$} \mtag{t}
\end{multicols}

% =================================================================== 3
\section{Equations with a square root}\label{h4:surd}
Square both sides, solve the quadratic, then \textbf{check every answer in the original equation}: squaring can create an answer that does not work, because $\sqrt{\ \ }$ is never negative. And $(x+2)^{2} = x^{2} + 4x + 4$, not $x^{2} + 4$.

\textbf{Example 7.} Solve $\sqrt{3x + 10} = x + 2$.
\[ 3x + 10 = x^{2} + 4x + 4,\qquad x^{2} + x - 6 = (x+3)(x-2) = 0 \]
Check $x = 2$: $\sqrt{16} = 4 = 2 + 2$. \quad Check $x = -3$: $\sqrt{1} = 1$ but $-3 + 2 = -1$, \textbf{rejected}. So $\mathbf{x = 2}$. \mtag{r}

% =================================================================== 4
\section{Modelling and constraints}\label{h4:model}
\begin{multicols}{2}
\textbf{Example 8.} A garden is 3 m longer than it is wide, with area 70 m$^{2}$. Find its size.
\[ x(x+3) = 70,\quad (x+10)(x-7) = 0 \]
A width cannot be negative, so $x = 7$: \textbf{7 m by 10 m}. \mtag{u}
\columnbreak

\textbf{Example 9.} A rectangle has perimeter 40 m. What is the largest area it can have?\par
Sides $x$ and $20 - x$: $x(20 - x) = A$, so $x^{2} - 20x + A = 0$.\par
A real width needs $400 - 4A \geq 0$, so $A \leq 100$: \textbf{at most 100 m$^{2}$}, a 10 m square. \mtag{t}
\end{multicols}
This is the move behind ``find the largest value of $a$'': write the condition as a quadratic, then ask when it has real solutions, then check the solutions against the context (a length must be positive).

% =================================================================== 5
\section{Versions the marker won't accept}\label{h4:wontaccept}
\begin{tabularx}{\textwidth}{@{}>{\raggedright}p{0.3\textwidth} >{\raggedright}p{0.26\textwidth} X@{}}
\toprule
\textbf{Question} & \textbf{Answer given} & \textbf{Why it falls short} \\ \midrule
Prove the sums of the corners are always equal & three squares checked by hand & Examples are not a proof. Use $A$ for the corner and write every corner from it. \tabularnewline
Find the co-ordinates of each intersection point & $x = \pm\sqrt{5}$, $x = \pm\sqrt{2}$ & Co-ordinates need $y$ too: four points, $(\pm\sqrt{5}, 2)$ and $(\pm\sqrt{2}, -1)$. \tabularnewline
Largest whole-number $a$ & $a = 7$, because $a^{2} < 50$ & The context also needs $p - a > 0$: with $a = 7$ that fails, so $a = 6$. \tabularnewline
\bottomrule
\end{tabularx}

% =================================================================== 6
\section{Ways to lose marks}\label{h4:lose}
From the wording of the marking schedules and from teaching observation; numbered for marking.
\begin{description}[leftmargin=10mm, labelwidth=9mm, labelsep=1mm, font=\normalfont]
\item[\textbf{\Lref{squaresum}}] \textbf{\Ltext{squaresum}.} $(x+3)^{2} = x^{2} + 6x + 9$.
\item[\textbf{\Lref{check}}] \textbf{\Ltext{check}.} Put each answer back into the square-root equation.
\item[\textbf{\Lref{general}}] \textbf{\Ltext{general}.} Letters for every case, then a sentence.
\item[\textbf{\Lref{allpoints}}] \textbf{\Ltext{allpoints}.} $x^{2} = 4$ gives $x = 2$ \textbf{and} $x = -2$; each $x$ needs its $y$.
\item[\textbf{\Lref{contextcheck}}] \textbf{\Ltext{contextcheck}.} Lengths and speeds are positive; sides like $p - a$ must be too.
\end{description}

% =================================================================== 7
\section{Quick review}\label{h4:review}
\renewcommand\arraystretch{1.45}
\begin{minipage}[t]{0.47\textwidth}\vspace{0pt}
\textbf{The discriminant (Lesson 3)}\par\smallskip
\begin{tabular}{@{}l l@{}}
\toprule
$b^{2} - 4ac > 0$ & two different real roots \\
$b^{2} - 4ac = 0$ & one repeated root (tangent) \\
$b^{2} - 4ac < 0$ & no real roots (no meeting) \\
\bottomrule
\end{tabular}
\end{minipage}\hfill
\begin{minipage}[t]{0.47\textwidth}\vspace{0pt}
\textbf{Squares to expand by sight}\par\smallskip
\begin{tabular}{@{}l l@{}}
\toprule
$(x + 2)^{2}$ & $x^{2} + 4x + 4$ \\
$(x - 5)^{2}$ & $x^{2} - 10x + 25$ \\
$(2n + 1)^{2}$ & $4n^{2} + 4n + 1$ \\
\bottomrule
\end{tabular}
\end{minipage}
\end{document}
